%! Function Model Selection and Assumption Articulation — Unit 1, Lesson 1.13 — Student Packet Cover Sheet
\documentclass[10pt]{article}
\usepackage{precalculus-article}
\usepackage{precalculus-boxes}
\usepackage{ltablex}
\keepXColumns

\begin{document}

% Full-bleed navy banner
\begin{tikzpicture}[remember picture, overlay]
  \fill[navy] (current page.north west)
    rectangle ([yshift=-0.9in]current page.north east);
\end{tikzpicture}
\vspace{-0.8in}

\begin{center}
  {\color{white}\textbf{\LARGE AP Precalculus}}\\[4pt]
  {\color{skymid}\large\textbf{Unit 1: Polynomial and Rational Functions}}\\[6pt]
  {\color{white}\textbf{\large Lesson 1.13 \quad Function Model Selection and Assumption Articulation}}
\end{center}

\vspace{0.22in}
\namedateperiod
\vspace{0.18in}

\begin{learningtargetbox}
\small
\begin{enumerate}[leftmargin=1.5em, itemsep=5pt]
  \item I can identify the appropriate function type (linear, quadratic, cubic, polynomial,
        piecewise) for a given data set or contextual scenario by analyzing rates of change.
  \item I can articulate the assumptions and domain/range restrictions embedded in a
        function model.
\end{enumerate}
\end{learningtargetbox}

\vspace{0.14in}

\begin{tocbox}
\small
\renewcommand{\arraystretch}{1.5}
\begin{tabularx}{\linewidth}{@{} c l X r @{}}
\rowcolor{sky}
\textbf{\#} & \textbf{Component} & \textbf{Description} & \textbf{Score} \\
\midrule
1 & Warm-Up        & Spiral review: first and second differences, constant vs.\ changing rates & \blank{1.2cm} \\
2 & Guided Notes   & Difference-table decision tree; shape/context clues; assumption templates & \blank{1.2cm} \\
3 & Group Activity & Data Dispatch Agency --- select models and state assumptions for scenarios & \blank{1.2cm} \\
4 & Exit Ticket    & Difference table, model classification, and assumption statement & \blank{1.2cm} \\
5 & Homework       & Independent practice: differences, matching, assumption articulation & \blank{1.2cm} \\
\midrule
  & \multicolumn{2}{r}{\textbf{Total}} & \blank{1.2cm} \\
\end{tabularx}
\end{tocbox}

\end{document}
